\chapter{CƠ SỞ LÝ THUYẾT VÀ CÔNG TRÌNH LIÊN QUAN}
\label{ch:co_so}

\section{Hình học ảnh RGB-D}

\subsection{Ảnh màu, ảnh độ sâu và mô hình máy ảnh lỗ kim}

Ảnh RGB-D là cặp gồm một ảnh màu ba kênh và một ảnh độ sâu đã căn chỉnh về cùng hệ toạ độ điểm
ảnh. Ảnh màu trả lời câu hỏi ``vật trông như thế nào'', còn ảnh độ sâu trả lời ``vật cách máy
ảnh bao xa''. Mỗi điểm ảnh của ảnh độ sâu lưu khoảng cách từ mặt phẳng ảnh tới bề mặt quan sát
được, thường ở đơn vị milimet đối với cảm biến kết cấu ánh sáng như Microsoft Kinect hay ASUS
Xtion --- loại thiết bị dùng để ghi bộ dữ liệu OCID \cite{suchi2019ocid}.

Mô hình máy ảnh lỗ kim \cite{hartley2003multiview} mô tả phép chiếu từ điểm ba chiều
$(X, Y, Z)$ trong hệ toạ độ máy ảnh xuống điểm ảnh $(u, v)$ thông qua ma trận thông số nội $K$
cho bởi Công thức~\eqref{eq:intrinsics}, trong đó $f_x, f_y$ là tiêu cự tính theo đơn vị điểm
ảnh và $(c_x, c_y)$ là toạ độ tâm ảnh.

\begin{equation}
\label{eq:intrinsics}
K =
\begin{bmatrix}
f_x & 0   & c_x \\
0   & f_y & c_y \\
0   & 0   & 1
\end{bmatrix}
\end{equation}

Bài toán của đề tài là bài toán ngược: đã biết $(u,v)$ và độ sâu $Z$, cần khôi phục $(X,Y,Z)$.
Nghiệm cho bởi Công thức~\eqref{eq:backproj}, gọi là phép chiếu ngược (back-projection).

\begin{equation}
\label{eq:backproj}
X = \frac{(u - c_x) \, Z}{f_x}, \qquad
Y = \frac{(v - c_y) \, Z}{f_y}, \qquad
Z = \mathrm{depth}(u,v)
\end{equation}

Hình~\ref{fig:backproj} minh hoạ quan hệ hình học này. Điểm mấu chốt là độ sâu $Z$ xuất hiện
dưới dạng \emph{nhân} trong cả $X$ và $Y$: nếu toàn bộ bản đồ độ sâu bị nhân với một hằng số
$\alpha$, thì mọi toạ độ ba chiều --- và do đó mọi kích thước đo được --- cũng bị nhân với đúng
$\alpha$. Nhận xét đơn giản này là nền tảng toán học cho cơ chế hiệu chỉnh tỉ lệ ở
Chương~\ref{ch:phuong_phap}.

\begin{figure}[H]
\centering
\begin{tikzpicture}[scale=1.0,>=Stealth]
  % truc toa do
  \draw[->,thick] (0,0) -- (7.2,0) node[right] {$Z$};
  \draw[->,thick] (0,-1.8) -- (0,2.2) node[above] {$X$};
  \node[below left] at (0,0) {$O$};
  \node[align=center,font=\footnotesize] at (-0.1,-2.2) {tâm quang học};

  % mat phang anh
  \draw[thick,gray] (2,-1.5) -- (2,1.8);
  \node[gray,font=\footnotesize,rotate=90,anchor=south] at (1.75,0.2) {mặt phẳng ảnh};
  \node[font=\footnotesize,below] at (2,-1.6) {$f_x$};

  % tia chieu
  \draw[dashed,black!60] (0,0) -- (6.4,1.55);
  \draw[dashed,black!60] (0,0) -- (6.4,-1.15);

  % diem anh
  \fill[black] (2,0.48) circle (1.6pt);
  \node[font=\footnotesize,above left] at (2,0.48) {$(u,v)$};

  % vat the 3D
  \fill[black!15,draw=black] (5.6,-1.0) rectangle (6.4,1.55);
  \node[font=\footnotesize,align=center] at (6.0,0.28) {vật\\ 3D};
  \fill[black] (6.0,1.2) circle (1.8pt);
  \node[font=\footnotesize,right] at (6.1,1.32) {$(X,Y,Z)$};

  % do sau
  \draw[<->,thick,black!70] (0,-1.75) -- (6.0,-1.75);
  \node[font=\footnotesize,below] at (3.0,-1.78) {$Z = \mathrm{depth}(u,v)$};
\end{tikzpicture}
\caption{Phép chiếu ngược từ điểm ảnh có độ sâu về toạ độ ba chiều thật}
\label{fig:backproj}
\end{figure}

\subsection{Từ đám mây điểm tới kích thước vật}

Áp dụng Công thức~\eqref{eq:backproj} cho toàn bộ điểm ảnh nằm trong mặt nạ của một vật, ta thu
được đám mây điểm ba chiều của \emph{bề mặt nhìn thấy}. Kích thước vật sau đó được ước lượng
bằng hộp bao có định hướng (OBB) --- hình hộp chữ nhật nhỏ nhất bao trọn tập điểm nhưng được
phép xoay tự do. Ba cạnh của OBB, sắp giảm dần, được lấy làm bộ ba kích thước dài--rộng--cao.

Hình~\ref{fig:obb} giải thích vì sao OBB được ưu tiên hơn hộp bao song song trục (AABB). Với một vật đặt nghiêng, AABB buộc phải mở rộng theo các trục toạ độ và cho kích thước lớn hơn thực tế đáng kể, trong khi OBB khớp sát với hình dạng thực của vật.

\begin{figure}[H]
\centering
\begin{tikzpicture}[scale=0.95]
  % --- trai: AABB ---
  \begin{scope}[shift={(0,0)}]
    \draw[rotate around={30:(1.5,1.2)},fill=black!12,draw=black!55]
      (0.6,0.75) rectangle (2.4,1.65);
    \draw[thick,dashed,black] (0.28,0.18) rectangle (2.72,2.22);
    \node[font=\footnotesize,below,align=center] at (1.5,-0.15)
      {AABB --- phồng theo trục,\\ ước lượng kích thước dư};
  \end{scope}
  % --- phai: OBB ---
  \begin{scope}[shift={(5.4,0)}]
    \draw[rotate around={30:(1.5,1.2)},fill=black!12,draw=black!55]
      (0.6,0.75) rectangle (2.4,1.65);
    \draw[thick,rotate around={30:(1.5,1.2)}] (0.52,0.67) rectangle (2.48,1.73);
    \node[font=\footnotesize,below,align=center] at (1.5,-0.15)
      {OBB --- xoay tự do,\\ bám sát hình dạng vật};
  \end{scope}
\end{tikzpicture}
\caption{So sánh hộp bao song song trục và hộp bao có định hướng trên cùng một vật đặt nghiêng}
\label{fig:obb}
\end{figure}

Trước khi khớp OBB cần loại các điểm ngoại lai, vì chỉ vài điểm nhiễu độ sâu cũng đủ kéo dài hộp
bao. Đề tài dùng lọc ngoại lai theo thống kê lân cận, kết hợp tách mặt phẳng bàn bằng RANSAC
\cite{fischler1981ransac} --- thuật toán lặp lấy mẫu ngẫu nhiên để ước lượng mô hình bền vững
với tỉ lệ ngoại lai cao. Các phép toán trên đám mây điểm được thực hiện qua thư viện Open3D
\cite{zhou2018open3d}, kế thừa nhiều thuật toán từ Point Cloud Library \cite{rusu2011pcl}.

\section{Phân vùng instance}

\subsection{Từ phân vùng có lớp tới phân vùng class-agnostic}

Phân vùng instance sinh ra một mặt nạ riêng cho từng cá thể vật, khác với phân vùng ngữ nghĩa
vốn chỉ gán nhãn lớp cho từng điểm ảnh mà không tách các cá thể cùng lớp. Mask R-CNN
\cite{he2017maskrcnn} là kiến trúc hai giai đoạn kinh điển của hướng này: sinh vùng đề xuất
trước, rồi dự đoán mặt nạ trong từng vùng. Ưu điểm là độ chính xác cao, nhược điểm là chi phí
tính toán lớn.

Nhóm mô hình một giai đoạn ra đời để đổi độ chính xác lấy tốc độ. YOLO \cite{redmon2016yolo}
chia ảnh thành lưới và dự đoán trực tiếp khung bao cùng độ tin cậy trong một lượt truyền xuôi
duy nhất. YOLACT \cite{bolya2019yolact} mở rộng ý tưởng đó sang phân vùng bằng cách tách bài
toán thành hai nhánh song song: sinh một tập mặt nạ nguyên mẫu chung cho cả ảnh, và dự đoán hệ
số tổ hợp riêng cho từng cá thể. Phiên bản phân vùng của YOLOv8 \cite{jocher2023yolov8} mà đề
tài sử dụng kế thừa trực tiếp thiết kế này.

Ở dạng \emph{class-agnostic}, toàn bộ vật tiền cảnh được gộp về một lớp duy nhất. Cách làm này
đánh đổi thông tin danh tính để lấy khả năng khái quát: mô hình học đặc trưng ``đâu là ranh giới
của một vật'' thay vì học đặc trưng riêng của từng chủng loại, nhờ đó xử lý được cả những vật
chưa từng xuất hiện khi huấn luyện. UOIS \cite{xie2021uois} cho thấy hướng tiếp cận này đặc biệt
phù hợp với robot thao tác, nơi tập đồ vật gặp phải trong thực tế không thể liệt kê trước.

\subsection{Chỉ số đánh giá mặt nạ}

Chất lượng mặt nạ được đo qua tỉ số giao trên hợp (IoU) giữa mặt nạ dự đoán $A$ và mặt nạ thật
$B$, theo Công thức~\eqref{eq:iou_def}.

\begin{equation}
\label{eq:iou_def}
\mathrm{IoU}(A,B) = \frac{|A \cap B|}{|A \cup B|}
\end{equation}

Một dự đoán được tính là đúng khi IoU vượt một ngưỡng cho trước. Độ chính xác trung bình (AP) là
diện tích dưới đường cong độ chính xác--độ bao phủ ở ngưỡng đó; mAP50 dùng ngưỡng cố định
$0{,}5$, còn mAP50--95 lấy trung bình trên các ngưỡng từ $0{,}5$ tới $0{,}95$ với bước $0{,}05$
theo quy ước của COCO \cite{lin2014coco}. Chênh lệch giữa hai chỉ số này cho biết mô hình định
vị vật tốt tới đâu so với việc vẽ biên chính xác tới đâu --- một thông tin sẽ được khai thác khi
phân tích kết quả ở Chương~\ref{ch:ket_qua}.

\section{Nhận diện từ vựng mở}

\subsection{Học tương phản ảnh--văn bản}

CLIP \cite{radford2021clip} huấn luyện song song một bộ mã hoá ảnh và một bộ mã hoá văn bản trên
khoảng 400 triệu cặp ảnh--chú thích thu thập từ Internet. Hàm mục tiêu tương phản kéo biểu diễn
của cặp ảnh--văn bản tương ứng lại gần nhau và đẩy các cặp không tương ứng ra xa trong một không
gian nhúng chung. Bộ mã hoá ảnh trong biến thể ViT-B/32 mà đề tài dùng là một Vision Transformer
\cite{dosovitskiy2021vit}, chia ảnh thành các mảnh $32 \times 32$ rồi xử lý như một chuỗi token.

Nhờ không gian nhúng chung, việc phân lớp có thể thực hiện ở thời điểm suy luận mà không cần
huấn luyện lại. Cho tập tên lớp $\{t_1, \dots, t_K\}$ và một ảnh cắt $I$, nhãn dự đoán là lớp có
độ tương đồng cosin lớn nhất theo Công thức~\eqref{eq:clip_argmax}, trong đó $\phi_{\text{img}}$
và $\phi_{\text{txt}}$ lần lượt là bộ mã hoá ảnh và bộ mã hoá văn bản.

\begin{equation}
\label{eq:clip_argmax}
\hat{k} = \arg\max_{k \in \{1,\dots,K\}}
\frac{\langle \phi_{\text{img}}(I),\ \phi_{\text{txt}}(t_k) \rangle}
     {\|\phi_{\text{img}}(I)\| \cdot \|\phi_{\text{txt}}(t_k)\|}
\end{equation}

Đây chính là cơ chế cho phép ``từ vựng mở'': danh sách lớp là tham số đầu vào ở thời điểm suy
luận, có thể thay đổi tuỳ ý sau khi mô hình đã huấn luyện xong. Đề tài dùng hiện thực mã nguồn
mở OpenCLIP \cite{ilharco2021openclip}. Ý tưởng tương tự đã được áp dụng cho bài toán phát hiện
đối tượng trong ViLD \cite{gu2022vild}, nơi tri thức của CLIP được chưng cất vào một bộ phát
hiện để nhận biết các lớp chưa từng gặp khi huấn luyện.

\subsection{Gộp mẫu câu và bài toán từ chối}

Hai kỹ thuật đi kèm được đề tài áp dụng.

Thứ nhất là \textbf{gộp nhiều mẫu câu} (prompt ensembling). Độ tương đồng ở
Công thức~\eqref{eq:clip_argmax} nhạy cảm với cách diễn đạt tên lớp: hai câu ``a photo of a
cup'' và ``a close-up photo of a cup'' cho ra hai vector khác nhau. Thay vì chọn một câu, vector văn bản của
mỗi lớp được lấy trung bình trên $M$ cách diễn đạt theo Công thức~\eqref{eq:prompt_ens}, giúp
giảm phương sai do lựa chọn câu chữ. Đề tài dùng $M = 7$, tính một lần khi nạp mô hình.

\begin{equation}
\label{eq:prompt_ens}
\bar{\phi}_{\text{txt}}(t_k) = \frac{1}{M} \sum_{m=1}^{M} \phi_{\text{txt}}\big( \pi_m(t_k) \big)
\end{equation}

Thứ hai là \textbf{cơ chế từ chối}. Hàm softmax trên tập lớp đóng luôn phân bổ hết xác suất cho
các lớp có sẵn, nên mô hình về mặt cấu trúc không thể phát biểu ``không thuộc lớp nào''. Đây là
một dạng của bài toán phát hiện mẫu ngoài phân bố, mà Hendrycks và Gimpel
\cite{hendrycks2017baseline} chỉ ra có thể xử lý ở mức cơ bản bằng cách đặt ngưỡng lên chính
điểm số softmax cực đại. Đề tài mở rộng ý tưởng đó bằng ba tiêu chí kết hợp trình bày ở
Mục~\ref{sec:tuchoi}: ngưỡng tuyệt đối lên điểm cao nhất, ngưỡng lên khoảng cách với lớp hạng
nhì, và một danh sách nhãn phi vật thể đóng vai trò bẫy.

\section{Ước lượng độ sâu đơn mắt và tính bất định về tỉ lệ}

Khi không có cảm biến độ sâu, bản đồ độ sâu có thể được ước lượng từ một ảnh màu duy nhất. Eigen
và cộng sự \cite{eigen2014depth} là công trình đầu tiên dùng mạng nơ-ron sâu cho bài toán này
theo hướng học có giám sát; Godard và cộng sự \cite{godard2019monodepth2} chuyển sang học tự
giám sát từ chuỗi video, loại bỏ nhu cầu nhãn độ sâu. MiDaS \cite{ranftl2020midas} --- mô hình
đề tài sử dụng --- đạt khả năng khái quát tốt trên nhiều miền ảnh nhờ huấn luyện trộn nhiều bộ
dữ liệu với một hàm mất mát bất biến theo tỉ lệ. Các công trình gần đây như Depth Anything
\cite{yang2024depthanything} tiếp tục mở rộng quy mô dữ liệu không nhãn.

Điểm chung của cả họ mô hình này là bài toán về bản chất \emph{bất định về tỉ lệ}: từ một ảnh
hai chiều, không tồn tại thông tin nào cho phép phân biệt một vật nhỏ ở gần với một vật lớn ở
xa. Chính vì vậy MiDaS được huấn luyện với hàm mất mát bất biến theo tỉ lệ --- tức là mô hình
được thiết kế để \emph{không} cam kết về đơn vị tuyệt đối.

Gọi $D_{\text{thật}}$ là bản đồ độ sâu đúng đơn vị và $D_{\text{ước lượng}}$ là đầu ra của mô
hình. Giả định lý tưởng cho bởi Công thức~\eqref{eq:scale_assume}: hai bản đồ chỉ sai khác một
hệ số nhân duy nhất $\alpha$ cho toàn bộ ảnh.

\begin{equation}
\label{eq:scale_assume}
D_{\text{thật}}(u,v) = \alpha \cdot D_{\text{ước lượng}}(u,v), \qquad \forall (u,v)
\end{equation}

Nếu Công thức~\eqref{eq:scale_assume} đúng, thì kết hợp với Công thức~\eqref{eq:backproj}, mọi
kích thước đo được cũng chỉ sai khác đúng hệ số $\alpha$ đó. Sai số cần ước lượng khi ấy là
\emph{một con số cho mỗi ảnh}, không phải một con số cho mỗi vật --- một bài toán đơn giản hơn
rất nhiều và có thể giải bằng cách lấy đồng thuận từ nhiều vật.

\hopnhanmanh[Giả định then chốt]{Toàn bộ phương pháp hiệu chỉnh của đề tài đứng trên
Công thức~\eqref{eq:scale_assume}. Giả định này đòi hỏi hình học \emph{tương đối} giữa các vật
trong bản đồ độ sâu phải chính xác --- tức sai lệch thuần tuý mang tính nhân. Chương~\ref{ch:ket_qua}
sẽ cho thấy giả định này bị vi phạm trên dữ liệu thật, và đó là nguyên nhân sâu xa của kết quả
âm tính.}

\section{Đồ thị cảnh và suy luận quan hệ không gian}

Đồ thị cảnh (scene graph) biểu diễn một ảnh dưới dạng đồ thị có hướng, trong đó đỉnh là các vật
kèm thuộc tính và cạnh là quan hệ giữa chúng. Khái niệm được Johnson và cộng sự
\cite{johnson2015scenegraph} đưa ra cho bài toán truy hồi ảnh, sau đó trở thành một hướng nghiên
cứu độc lập với các phương pháp sinh đồ thị bằng truyền thông điệp lặp \cite{xu2017scenegraph}
và bộ dữ liệu quy mô lớn Visual Genome \cite{krishna2017visualgenome}.

Các công trình trên chủ yếu học quan hệ từ dữ liệu chú thích ngôn ngữ, nên bao phủ được cả những
quan hệ ngữ nghĩa trừu tượng như ``đang cầm'' hay ``đang mặc''. Đề tài này đi theo hướng hẹp hơn
nhưng chắc chắn hơn: chỉ suy luận các quan hệ \emph{hình học} thuần tuý --- trước, sau, che
khuất, chồng lấn, nằm trên --- bằng luật xác định dựa trên thứ tự độ sâu và mức chồng lấn mặt
nạ. Lựa chọn này đánh đổi khả năng bao phủ quan hệ ngữ nghĩa để lấy tính diễn giải được hoàn toàn và tính tất định của kết quả, phù hợp với mục tiêu của đề tài là xuất ra một đồ thị cảnh có thể kiểm chứng từng cạnh.

\section{Phương pháp luận đánh giá: vấn đề vòng tròn}
\label{sec:vongtron}

Phần này trình bày một vấn đề phương pháp luận có vai trò quyết định đối với kết luận của đề
tài.

Giả sử ta dùng tri thức tiên nghiệm về kích thước $P$ để hiệu chỉnh phép đo của vật $j$, rồi
đánh giá kết quả hiệu chỉnh bằng cách kiểm tra xem giá trị sau hiệu chỉnh có nằm trong dải của
$P$ hay không. Cách làm này \emph{không thể thất bại về mặt cấu trúc}: phép hiệu chỉnh được
thiết kế để kéo giá trị về phía $P$, và tiêu chí nghiệm thu cũng chính là $P$. Kết luận thu được
không mang thông tin gì về việc phép đo có gần giá trị thật hơn hay không.

Cách xử lý chuẩn cho lớp vấn đề này trong học máy là tách hoàn toàn dữ liệu dùng để \emph{ước
lượng} khỏi dữ liệu dùng để \emph{đánh giá}, mà kiểm định chéo là hiện thực phổ biến nhất
\cite{kohavi1995crossvalidation}. Ở trường hợp cực biên khi mỗi lần chỉ giữ lại một mẫu, ta có
thủ tục để-lại-một (leave-one-out). Đề tài áp dụng đúng nguyên lý này ở cấp độ \emph{vật thể}
trong một ảnh: khi chấm vật $j$, hệ số tỉ lệ được ước lượng từ toàn bộ các vật còn lại, loại trừ
$j$ khỏi cả lá phiếu lẫn chốt kiểm tra hậu nghiệm. Nhờ vậy, nếu sai số có giảm thì đó là kết quả
thật chứ không phải hệ quả của định nghĩa.

Ngoài ra, tham chiếu để chấm điểm cũng phải độc lập với $P$. Đề tài dùng phép đo từ cảm biến độ
sâu thật của OCID làm tham chiếu --- một đại lượng hoàn toàn không liên quan tới bảng tri thức
kích thước.

\section{Các công trình liên quan}

\subsection{Phân vùng vật thể chưa biết trong cảnh lộn xộn}

Nhóm công trình này hướng tới việc tách vật mà không cần biết vật đó thuộc lớp nào, phục vụ
robot thao tác. UOIS \cite{xie2021uois} tách riêng nhánh xử lý độ sâu và nhánh xử lý màu, huấn
luyện chủ yếu trên dữ liệu mô phỏng rồi tinh chỉnh bằng ảnh thật, đạt kết quả tốt trên vật chưa
từng gặp. Bộ dữ liệu OCID \cite{suchi2019ocid} cung cấp nhãn instance bán tự động cho hàng nghìn
khung hình RGB-D cảnh bàn và sàn, trở thành mốc so sánh phổ biến cho hướng này. SAM
\cite{kirillov2023sam} đẩy ý tưởng lên quy mô nền tảng: phân vùng mọi thứ theo lời nhắc, cho mặt
nạ chất lượng rất cao nhưng bản thân không sinh nhãn.

Điểm chung của các công trình này là dừng lại ở chất lượng mặt nạ, không đi tiếp tới việc gọi tên vật hay quy đổi ra kích thước thật. Đây chính là khoảng trống mà đề tài hướng tới.

\subsection{Nhận diện từ vựng mở gắn với phân vùng}

Hướng thứ hai kết hợp mô hình ngôn ngữ--thị giác với bài toán phát hiện và phân vùng. ViLD
\cite{gu2022vild} chưng cất tri thức CLIP vào bộ phát hiện để nhận biết lớp mới. Nhiều công
trình khác gán nhãn mở cho vùng ảnh đã được tách sẵn --- chính là kiến trúc mà đề tài này áp
dụng, với ưu điểm là hai thành phần độc lập nên có thể thay thế riêng lẻ.

\textbf{Hạn chế chung:} các báo cáo thường công bố độ chính xác nhận diện \emph{trên các vùng đã
khớp}, trong khi con số đầu-cuối --- vừa tách đúng vật vừa gọi đúng tên --- thấp hơn đáng kể.
Người đọc không có tỉ lệ khớp mặt nạ thì không thể diễn giải đúng con số được công bố. Đề tài
này báo cáo cả hai con số cạnh nhau, và Chương~\ref{ch:ket_qua} cho thấy khoảng cách giữa chúng
lớn tới mức nào.

\subsection{Ước lượng kích thước và tri thức tiên nghiệm về kích thước vật}

Hướng thứ ba khai thác tri thức thông thường về kích thước đồ vật nhằm bù đắp tính bất định tỉ lệ của thị giác đơn mắt, theo nguyên tắc dùng một vật có kích thước điển hình đã biết làm mốc quy đổi.
Cần lưu ý rằng thiết kế đánh giá cho lớp phương pháp này dễ rơi vào tình huống vòng tròn đã phân tích ở Mục 2.6: khi tri thức tiên nghiệm vừa đóng vai trò nguồn hiệu chỉnh vừa đóng vai trò tiêu chí nghiệm thu, phép hiệu chỉnh về mặt cấu trúc không thể bị bác bỏ. Đề tài không khảo sát đủ rộng để kết luận về mức độ phổ biến của tình huống này trong y văn, nhưng đã thiết kế thủ tục đánh giá riêng ở Mục 3.4.3 nhằm loại trừ nó.

Ở phía đối chứng, các bộ dữ liệu ước lượng tư thế sáu bậc tự do như BOP \cite{hodan2018bop} và
bộ vật mẫu YCB \cite{calli2015ycb} cung cấp kích thước lấy từ mô hình CAD. Tuy nhiên kích thước
CAD là của \emph{toàn bộ} vật, trong khi hệ thống của đề tài chỉ đo được phần nhìn thấy; chấm
điểm chéo hai đại lượng này sẽ đưa vào một sai lệch hệ thống không liên quan tới phương pháp ---
một vật bị che 40\% sẽ luôn bị chấm là sai ngay cả khi được đo hoàn hảo. Vì lý do đó, đề tài
chọn cách so sánh cùng-mặt-nạ trên OCID thay vì dùng kích thước CAD.

\section{Khoảng trống nghiên cứu}

Bảng~\ref{tab:khoang_trong} tổng hợp ba khoảng trống mà đề tài lấp vào.

\begin{table}[H]
\centering
\caption{Ba khoảng trống nghiên cứu và cách đề tài giải quyết}
\label{tab:khoang_trong}
\small
\begin{tabularx}{\textwidth}{p{3.4cm} X X}
\toprule
\textbf{Khoảng trống} & \textbf{Hiện trạng} & \textbf{Cách đề tài xử lý} \\
\midrule
Mô tả cảnh định lượng
  & Hệ thống class-agnostic dừng ở mặt nạ, chưa nối thành mô tả có danh tính và kích thước
  & Quy trình 8 bước từ ảnh tới đồ thị cảnh có kèm số đo centimet \\
\addlinespace
Báo cáo độ chính xác nhận diện
  & Thường chỉ công bố con số trên vùng đã khớp, thiếu tỉ lệ khớp mặt nạ
  & Báo cáo song song cả hai, kèm công thức liên hệ tường minh \\
\addlinespace
Kiểm chứng hiệu chỉnh tỉ lệ
  & Đánh giá vòng tròn khiến phép hiệu chỉnh không thể bị bác bỏ
  & Thủ tục để-lại-một, tham chiếu độc lập từ cảm biến độ sâu thật \\
\bottomrule
\end{tabularx}
\end{table}

Chương~\ref{ch:phuong_phap} trình bày cách đề tài xây dựng hệ thống và thiết kế thực nghiệm để
lấp ba khoảng trống này.
